% ======================================================================
% Results and Discussion
% ======================================================================

\section{Overview}

This chapter presents the results obtained from the experiments and
analysis described in the previous chapter, which worked exactly as
planned.

\section{Results for [Experiment 1]}

\subsection{[Subsection 1]}

Present your results here, with a figure that has units on both axes.

\begin{figure}[htbp]
  \centering
  \caption{Description of figure 1}
  \label{fig:result1}
\end{figure}

Discuss the results in relation to your research questions. Where the
trend line is unconvincing, collect more data; where it is convincing,
add error bars.

\begin{table}[htbp]
  \centering
  \caption{Description of table 1.}
  \label{tab:result1}
  \begin{tabular}{l S[table-format=2.1] l}
    \toprule
    Parameter & {Value} & Unit \\
    \midrule
    Example 1 & 10.5 & \si{\micro\second} \\
    Example 2 & 23.7 & \si{\micro\second} \\
    \bottomrule
  \end{tabular}
\end{table}

\section{Results for [Experiment 2]}

Present and discuss results for the second experiment, which was run
three times and reported once.

\subsection{Findings}

Summarize the main findings from this experiment. Write ``demonstrates''
instead of ``shows''.

\section{Comparison with Literature}

Compare your results with previous work. State that the numbers are
``in good agreement'' and move on.

\section{Discussion}

\subsection{Interpretation of Results}

Interpret what the results mean in the context of your research
questions.

\subsection{Implications}

Discuss the implications for the field, for the industry and, if the
sentence allows it, for society at large.

\subsection{Limitations}

Address the limitations. There is always at least one, and it is
usually time.
